\documentclass[12pt]{article}

\usepackage[utf8]{inputenc}
\usepackage[T1]{fontenc}
\usepackage{lmodern}
\usepackage[margin=1in]{geometry}
\usepackage{graphicx}
\usepackage{hyperref}
\usepackage{array}
\usepackage{booktabs}
\usepackage{enumitem}
\usepackage{listings}
\usepackage{xcolor}
\usepackage{float}

\setlist{nosep}

% Code listing styling
\definecolor{codegray}{rgb}{0.5,0.5,0.5}
\definecolor{codepurple}{rgb}{0.58,0,0.82}
\definecolor{backcolour}{rgb}{0.95,0.95,0.92}

\lstdefinestyle{pythonstyle}{
    backgroundcolor=\color{backcolour},
    commentstyle=\color{codegray},
    keywordstyle=\color{blue},
    numberstyle=\tiny\color{codegray},
    stringstyle=\color{codepurple},
    basicstyle=\ttfamily\footnotesize,
    breakatwhitespace=false,
    breaklines=true,
    captionpos=b,
    keepspaces=true,
    numbers=left,
    numbersep=5pt,
    showspaces=false,
    showstringspaces=false,
    showtabs=false,
    tabsize=2,
    language=Python
}

\lstset{style=pythonstyle}

\title{Road Anomaly Detection System\\Technical Implementation Document}
\author{
  Atem Ako \and
  Stapanavichet Long \and
  Nahaeli Brunder \and
  Amir Abdalla
}
\date{November 2025}

\begin{document}

\maketitle

\tableofcontents
\clearpage

\section{Executive Summary}

This document describes the technical implementation of the Road Anomaly Detection System using the \texttt{road-anomaly-detection} repository as the core ML inference engine. The system leverages YOLOv8-based deep learning models to detect potholes, cracks, and other road surface anomalies from dashcam video feeds and static images.

\medskip

The implementation consists of two primary components:

\begin{itemize}
    \item \textbf{Frontend}: Streamlit-based web dashboard for real-time visualization, model configuration, and result analysis.
    \item \textbf{Backend}: Python-based ML processing pipeline utilizing YOLOv8 models for object detection, OpenCV for image processing, and the Supervision library for annotation management.
\end{itemize}

\section{Repository Overview}

\subsection{Repository Structure}

The \texttt{road-anomaly-detection} repository provides a complete implementation of the ML Inference Service described in the architectural design document.

\begin{table}[H]
    \centering
    \caption{Key Repository Components}
    \label{tab:repo-structure}
    \begin{tabular}{p{0.35\textwidth}p{0.55\textwidth}}
        \toprule
        \textbf{Component} & \textbf{Description} \\
        \midrule
        \texttt{main.py} & Streamlit web application with live camera, video, and image processing modes \\
        \texttt{run.py} & Standalone CLI script for batch video/image processing \\
        \texttt{YOLOv8\_Small\_2nd\_Model.pt} & Pre-trained YOLOv8s model (89.5 MB) for crack and pothole detection \\
        \texttt{RoadDetectionModel/} & Custom-trained YOLOv8m model with enhanced road anomaly detection \\
        \texttt{dataset/} & Training, validation, and test datasets in YOLOv8 format \\
        \texttt{requirements.txt} & Python dependencies (ultralytics, supervision, streamlit, opencv-python) \\
        \bottomrule
    \end{tabular}
\end{table}

\subsection{Trained Models}

The system employs two complementary YOLOv8 models:

\begin{table}[H]
    \centering
    \caption{Detection Models Comparison}
    \label{tab:models}
    \begin{tabular}{p{0.12\textwidth}p{0.25\textwidth}p{0.28\textwidth}p{0.22\textwidth}}
        \toprule
        \textbf{Model} & \textbf{Architecture} & \textbf{Detected Classes} & \textbf{Performance (mAP@0.5)} \\
        \midrule
        M1 & YOLOv8m (52 MB) & Heavy-Vehicle, Light-Vehicle, Pedestrian, Crack, Crack-Severe, Pothole, Speed-Bump & 0.745 \\
        M2 & YOLOv8s (89.5 MB) & Longitudinal Crack, Transverse Crack, Alligator Crack, Potholes & 0.40 (est.) \\
        \bottomrule
    \end{tabular}
\end{table}

\textbf{Model Selection Strategy}:

\begin{itemize}
    \item \textbf{M1} (Custom-trained): Higher accuracy for diverse road anomalies; trained on 30,685 images over 120 epochs.
    \item \textbf{M2} (Pre-trained): Specialized for crack classification; faster inference on edge devices.
    \item Both models can run concurrently for ensemble detection with color-coded bounding boxes (M1: red, M2: blue).
\end{itemize}

\clearpage

\section{Frontend: Web Dashboard Implementation}

\subsection{Streamlit Application Architecture}

The \texttt{main.py} module implements a fully functional web dashboard that serves as the primary user interface for the detection system. The application provides three operational modes:

\begin{enumerate}
    \item \textbf{Image Mode}: Upload and process static images.
    \item \textbf{Video Mode}: Upload video files for frame-by-frame analysis with downloadable annotated output.
    \item \textbf{Live Camera Mode}: Real-time detection from webcam or USB camera feed using WebRTC.
\end{enumerate}

\subsection{Key Frontend Features}

\subsubsection*{Model Configuration Panel}

Users can dynamically configure detection parameters:

\begin{itemize}
    \item \textbf{Model Selection}: Toggle M1 and/or M2 models independently.
    \item \textbf{Confidence Thresholds}: Adjust per-model confidence levels (0.1--1.0) via interactive sliders.
    \item \textbf{Input Source Selection}: Radio buttons for mode switching.
\end{itemize}

\subsubsection*{Real-Time Video Processing}

The video processing workflow:

\begin{enumerate}
    \item User uploads video file (MP4, AVI, MOV, MKV).
    \item System extracts video properties (FPS, resolution, frame count).
    \item Each frame is processed through the detection pipeline.
    \item Progress bar displays processing status with frame counters.
    \item Annotated video is generated and made available for download.
    \item Session state management prevents redundant reprocessing.
\end{enumerate}

\begin{lstlisting}[caption={Video Processing Core Logic (Simplified)}, label={lst:video-processing}]
def handle_video_input(models, thresholds, status_placeholder):
    # Video capture initialization
    cap = cv2.VideoCapture(input_tmp_path)
    width = int(cap.get(cv2.CAP_PROP_FRAME_WIDTH))
    height = int(cap.get(cv2.CAP_PROP_FRAME_HEIGHT))
    fps = cap.get(cv2.CAP_PROP_FPS)
    
    # Process each frame
    while True:
        ret, frame = cap.read()
        if not ret: break
        
        out_frame = process_frame(frame, models, thresholds)
        writer.write(out_frame)
        
        # Update progress bar
        prog_bar.progress(frame_idx / total_frames)
\end{lstlisting}

\subsubsection*{Live Camera Integration}

Utilizes \texttt{streamlit-webrtc} for browser-based camera access:

\begin{itemize}
    \item WebRTC peer connection with STUN server configuration.
    \item Asynchronous frame processing to maintain real-time performance.
    \item Automatic frame resizing to 640px width for inference optimization.
    \item Bidirectional video stream (SENDRECV mode).
\end{itemize}

\subsection{Detection Result Visualization}

The \texttt{supervision} library provides professional annotation capabilities:

\begin{itemize}
    \item \textbf{Bounding Boxes}: Color-coded by model (RED for M1, BLUE for M2).
    \item \textbf{Labels}: Display class name and confidence score (e.g., ``M1:Pothole 0.87'').
    \item \textbf{Overlay Text}: Frame counters and detection statistics.
\end{itemize}

\clearpage

\section{Backend: ML Processing Pipeline}

\subsection{Detection Engine Architecture}

The backend implements a three-stage detection pipeline aligned with the ML Subsystem specification in the design document.

\subsubsection*{Stage 1: Pre-Processing}

\begin{lstlisting}[caption={Frame Preprocessing}, label={lst:preprocessing}]
def preprocess_frame(frame):
    # Resize to target dimensions (640x640 for YOLOv8)
    # Note: YOLO models handle normalization internally
    
    # Optional: Enhance contrast for low-light conditions
    if is_low_light(frame):
        frame = cv2.convertScaleAbs(frame, alpha=1.2, beta=10)
    
    return frame
\end{lstlisting}

\subsubsection*{Stage 2: Inference}

The YOLOv8 inference is handled by the \texttt{ultralytics} library:

\begin{lstlisting}[caption={Detection Inference Call}, label={lst:inference}]
def process_frame(frame, models, thresholds):
    annotated_frame = frame.copy()
    
    for model_name, (model, names_map, box_ann, label_ann) in models.items():
        # Run inference
        results = model.predict(
            frame, 
            conf=thresholds[model_name], 
            verbose=False
        )[0]
        
        # Convert to supervision format
        detections = sv.Detections.from_ultralytics(results)
        
        # Generate labels
        labels = [
            f"{MODEL_PREFIX[model_name]}:{names_map[cls_id]} {conf:.2f}"
            for cls_id, conf in zip(detections.class_id, 
                                   detections.confidence)
        ]
        
        # Annotate frame
        annotated_frame = box_ann.annotate(annotated_frame, detections)
        annotated_frame = label_ann.annotate(
            annotated_frame, detections, labels=labels
        )
    
    return annotated_frame
\end{lstlisting}

\subsubsection*{Stage 3: Post-Processing}

Post-processing operations include:

\begin{itemize}
    \item \textbf{Non-Maximum Suppression (NMS)}: Removes overlapping bounding boxes (handled internally by YOLOv8).
    \item \textbf{Confidence Filtering}: Eliminates low-confidence detections based on user-defined thresholds.
    \item \textbf{Severity Estimation}: Maps bounding box dimensions to severity levels (Critical, High, Medium, Low).
\end{itemize}

\subsection{Batch Processing Mode}

The \texttt{run.py} script provides a command-line interface for batch processing:

\begin{lstlisting}[caption={Batch Processing Configuration}, label={lst:batch-config}]
# Configuration
BEST_WEIGHTS_PATH = "YOLOv8_Small_2nd_Model.pt"
CONF_THRESHOLD = 0.35
MODE = "video"  # Options: "image", "video", "live"

INPUT_VIDEO_PATH = r"path\to\video.mp4"
OUTPUT_DIR = "inference_output"
CAMERA_INDEX = 0  # For live mode
\end{lstlisting}

This mode is optimized for:

\begin{itemize}
    \item Processing large video archives.
    \item Headless server environments without GUI.
    \item Integration with automated data pipelines.
\end{itemize}

\subsection{Performance Characteristics}

\begin{table}[H]
    \centering
    \caption{Processing Performance Metrics}
    \label{tab:performance}
    \begin{tabular}{lll}
        \toprule
        \textbf{Metric} & \textbf{M1 (YOLOv8m)} & \textbf{M2 (YOLOv8s)} \\
        \midrule
        Inference Time (per image) & $\sim$12.0 ms & $\sim$8.0 ms (est.) \\
        Video Processing (30 FPS) & $\sim$2.5x real-time & $\sim$3.7x real-time \\
        Model Size & 52 MB & 89.5 MB \\
        GPU Memory Usage & $\sim$2.5 GB & $\sim$1.8 GB \\
        \bottomrule
    \end{tabular}
\end{table}

\textit{Note: Measurements taken on NVIDIA GeForce RTX 3060 Laptop GPU (6GB VRAM).}

\clearpage

\section{Hardware Integration Design}

\subsection{Edge Device Architecture}

The backend is designed to run on vehicle-mounted edge computing devices with the following specifications:

\begin{table}[H]
    \centering
    \caption{Recommended Edge Device Specifications}
    \label{tab:edge-specs}
    \begin{tabular}{ll}
        \toprule
        \textbf{Component} & \textbf{Specification} \\
        \midrule
        Processor & ARM Cortex-A72 or x86-64 (4+ cores) \\
        RAM & 4 GB minimum, 8 GB recommended \\
        GPU & NVIDIA Jetson Nano / Xavier NX (optional, for acceleration) \\
        Storage & 32 GB eMMC or SSD \\
        Camera Interface & USB 3.0 or CSI-2 (for Pi Camera modules) \\
        GPS Module & U-blox NEO-6M or similar (UART/I2C interface) \\
        Network & 4G/LTE modem or Wi-Fi with offline buffering \\
        Power & 12V DC (vehicle power supply compatible) \\
        \bottomrule
    \end{tabular}
\end{table}

\subsection{Data Acquisition Subsystem}

\subsubsection*{Camera Module}

The system supports multiple camera configurations:

\begin{itemize}
    \item \textbf{USB Webcams}: Standard UVC-compatible cameras (1080p @ 30 FPS recommended).
    \item \textbf{CSI Cameras}: Raspberry Pi Camera Module v2 or IMX219-based sensors.
    \item \textbf{IP Cameras}: Network cameras via RTSP streams (for cloud-connected deployments).
\end{itemize}

Frame capture logic uses OpenCV's \texttt{VideoCapture} API:

\begin{lstlisting}[caption={Camera Frame Capture}, label={lst:camera-capture}]
# Initialize camera
cap = cv2.VideoCapture(CAMERA_INDEX)
cap.set(cv2.CAP_PROP_FRAME_WIDTH, 1920)
cap.set(cv2.CAP_PROP_FRAME_HEIGHT, 1080)
cap.set(cv2.CAP_PROP_FPS, 30)

# Capture loop
while True:
    ret, frame = cap.read()
    if not ret:
        logger.error("Frame capture failed")
        continue
    
    # Process frame
    annotated = process_frame(frame, models, thresholds)
    
    # Display or store
    cv2.imshow("Detection Feed", annotated)
\end{lstlisting}

\subsubsection*{GPS Integration (Planned)}

\textbf{Current Status}: GPS hardware interface not yet implemented.

\medskip

\textbf{Proposed Implementation}:

\begin{enumerate}
    \item \textbf{GPS Module}: U-blox NEO-6M connected via USB-to-Serial adapter.
    \item \textbf{Python Library}: \texttt{gpsd} or \texttt{pynmea2} for NMEA sentence parsing.
    \item \textbf{Data Synchronization}: Timestamp alignment between GPS fixes and video frames using NTP-synchronized system clocks.
\end{enumerate}

\begin{lstlisting}[caption={GPS Data Acquisition (Conceptual)}, label={lst:gps-acquisition}]
import serial
import pynmea2

def read_gps_location():
    """Read GPS coordinates from serial port."""
    ser = serial.Serial('/dev/ttyUSB0', 9600, timeout=1)
    
    while True:
        line = ser.readline().decode('ascii', errors='ignore')
        
        if line.startswith('$GPGGA'):
            msg = pynmea2.parse(line)
            return {
                'latitude': msg.latitude,
                'longitude': msg.longitude,
                'altitude': msg.altitude,
                'timestamp': msg.timestamp,
                'fix_quality': msg.gps_qual
            }
\end{lstlisting}

\textbf{Integration with Detection Pipeline}:

Each detection would be augmented with GPS metadata:

\begin{lstlisting}[caption={Detection with GPS Metadata}, label={lst:detection-gps}]
class PotholeDetection:
    def __init__(self, bbox, confidence, severity, gps_data):
        self.x1, self.y1, self.x2, self.y2 = bbox
        self.confidence = confidence
        self.severity = severity
        
        # GPS metadata
        self.latitude = gps_data['latitude']
        self.longitude = gps_data['longitude']
        self.timestamp = gps_data['timestamp']
        self.fix_quality = gps_data['fix_quality']
\end{lstlisting}

\subsection{Data Upload \& Synchronization}

\subsubsection*{Offline Buffering}

The edge device implements a queuing mechanism for network-disconnected scenarios:

\begin{itemize}
    \item \textbf{Local SQLite Database}: Stores detection metadata and image references.
    \item \textbf{File System Buffer}: Compressed JPEG images stored in \texttt{/var/cache/pothole-buffer/}.
    \item \textbf{Upload Scheduler}: Background daemon monitors network connectivity and flushes queue when online.
\end{itemize}

\subsubsection*{Upload Protocol}

Detections are uploaded to the backend API via HTTP POST:

\begin{lstlisting}[caption={Detection Upload Request}, label={lst:upload-request}]
import requests
import base64

def upload_detection(detection, image_path):
    """Upload detection to backend API."""
    
    # Read and encode image
    with open(image_path, 'rb') as f:
        image_b64 = base64.b64encode(f.read()).decode('utf-8')
    
    payload = {
        'device_id': DEVICE_ID,
        'timestamp': detection.timestamp,
        'latitude': detection.latitude,
        'longitude': detection.longitude,
        'confidence': detection.confidence,
        'severity': detection.severity,
        'image_data': image_b64,
        'model_version': MODEL_VERSION
    }
    
    response = requests.post(
        f"{API_BASE_URL}/api/v1/detections",
        json=payload,
        headers={'Authorization': f'Bearer {API_TOKEN}'}
    )
    
    return response.status_code == 201
\end{lstlisting}

\subsection{Power Management}

To ensure continuous operation in vehicle environments:

\begin{itemize}
    \item \textbf{Deep Sleep Mode}: Enter low-power state when vehicle is off (accelerometer-based detection).
    \item \textbf{Voltage Monitoring}: Shutdown gracefully if vehicle battery drops below 11.5V.
    \item \textbf{Scheduled Processing}: Process frames only during daylight hours (6 AM--8 PM) to conserve power.
\end{itemize}

\clearpage

\section{Deployment \& Configuration}

\subsection{Installation Procedure}

\subsubsection*{Prerequisites}

\begin{itemize}
    \item Python 3.9 or higher
    \item pip package manager
    \item Git
    \item CUDA 11.8+ (optional, for GPU acceleration)
\end{itemize}

\subsubsection*{Setup Steps}

\begin{lstlisting}[language=bash, caption={Installation Commands}]
# Clone repository
git clone https://github.com/collabdoor/Road-Anomaly-Detection.git
cd Road-Anomaly-Detection

# Create virtual environment
python -m venv venv
source venv/bin/activate  # On Windows: venv\Scripts\activate

# Install PyTorch with CUDA support (if GPU available)
pip3 install torch torchvision torchaudio \
    --index-url https://download.pytorch.org/whl/cu128

# Install dependencies
pip install -r requirements.txt

# Verify installation
python -c "import torch; print(f'CUDA: {torch.cuda.is_available()}')"
\end{lstlisting}

\subsection{Running the System}

\subsubsection*{Web Dashboard}

\begin{lstlisting}[language=bash]
streamlit run main.py
\end{lstlisting}

The dashboard will be accessible at \texttt{http://localhost:8501}.

\subsubsection*{Batch Processing}

\begin{lstlisting}[language=bash]
# Edit run.py to configure input/output paths
python run.py
\end{lstlisting}

\subsubsection*{Live Camera Mode}

\begin{lstlisting}[language=bash]
# For standalone camera processing
python run.py  # Set MODE = "live" in configuration
\end{lstlisting}

\subsection{Configuration Parameters}

\begin{table}[H]
    \centering
    \caption{Key Configuration Parameters}
    \label{tab:config-params}
    \begin{tabular}{p{0.3\textwidth}p{0.5\textwidth}l}
        \toprule
        \textbf{Parameter} & \textbf{Description} & \textbf{Default} \\
        \midrule
        \texttt{MODEL\_PATHS} & Paths to trained model weights & See \texttt{main.py} \\
        \texttt{DEFAULT\_CONF} & Initial confidence thresholds & M1: 0.35, M2: 0.40 \\
        \texttt{LIVE\_FEED\_TARGET\_WIDTH} & Frame width for live processing & 640 px \\
        \texttt{OUTPUT\_DIR} & Directory for processed videos & \texttt{inference\_output/} \\
        \texttt{CAMERA\_INDEX} & Camera device index & 0 (default camera) \\
        \bottomrule
    \end{tabular}
\end{table}

\clearpage

\section{Integration with System Architecture}

\subsection{Mapping to Architectural Components}

The repository implementation directly maps to the following components from the design document:

\begin{table}[H]
    \centering
    \caption{Architecture-to-Implementation Mapping}
    \label{tab:arch-mapping}
    \begin{tabular}{p{0.35\textwidth}p{0.55\textwidth}}
        \toprule
        \textbf{Architectural Component} & \textbf{Repository Implementation} \\
        \midrule
        ML Inference Service & \texttt{process\_frame()} function with YOLOv8 model inference \\
        Data Ingestion Service & File upload handlers in \texttt{handle\_image\_input()} and \texttt{handle\_video\_input()} \\
        Storage Layer (Object Store) & Temporary file storage using \texttt{tempfile} module \\
        Dashboard / Reporting UI & Streamlit web application in \texttt{main.py} \\
        API Gateway & Streamlit's built-in HTTP server (not production-ready) \\
        Monitoring \& Logging & Python \texttt{logging} module with configurable levels \\
        \bottomrule
    \end{tabular}
\end{table}

\subsection{Missing Components (Future Work)}

The following architectural components from the design document are not yet implemented:

\begin{enumerate}
    \item \textbf{Pothole Management Service}: Business logic for merging duplicate detections, geo-clustering, and lifecycle status management.
    \item \textbf{Persistent Database}: No PostgreSQL or permanent storage; current implementation uses session state only.
    \item \textbf{Message Queue}: No asynchronous job processing via RabbitMQ/Kafka.
    \item \textbf{Authentication \& Authorization}: No user management or RBAC system.
    \item \textbf{GPS Hardware Integration}: GPS module interface not implemented.
    \item \textbf{Production API}: RESTful API endpoints not available (Streamlit is UI-only).
\end{enumerate}

\subsection{Proposed System Architecture}

\begin{figure}[H]
    \centering
    \fbox{\parbox{0.85\textwidth}{%
        \textbf{System Flow Diagram (Conceptual)}\\[10pt]
        \texttt{Edge Device (Vehicle)}\\
        \quad $\rightarrow$ Camera Module (\texttt{cv2.VideoCapture})\\
        \quad $\rightarrow$ GPS Module (\textit{planned}, U-blox NEO-6M)\\
        \quad $\rightarrow$ Detection Pipeline (\texttt{process\_frame()})\\
        \quad $\rightarrow$ Local Buffer (SQLite + File System)\\
        \quad $\rightarrow$ Upload Daemon\\[10pt]
        \texttt{Backend API Server}\\
        \quad $\leftarrow$ HTTP POST \texttt{/api/v1/detections}\\
        \quad $\rightarrow$ PostgreSQL Database\\
        \quad $\rightarrow$ Object Storage (S3/MinIO)\\[10pt]
        \texttt{Web Dashboard}\\
        \quad $\leftarrow$ API GET \texttt{/api/v1/potholes}\\
        \quad $\rightarrow$ Interactive Map (Leaflet/Mapbox)\\
        \quad $\rightarrow$ Filters \& Export
    }}
    \caption{End-to-End System Architecture with Repository Components}
    \label{fig:system-flow}
\end{figure}

\clearpage

\section{Testing \& Validation}

\subsection{Model Performance Validation}

\subsubsection*{Test Set Evaluation (Model M1)}

Comprehensive evaluation on 2,846 test images:

\begin{table}[H]
    \centering
    \caption{Model M1 Test Set Performance}
    \label{tab:model-test-performance}
    \begin{tabular}{lcccc}
        \toprule
        \textbf{Class} & \textbf{Precision} & \textbf{Recall} & \textbf{mAP@0.5} & \textbf{mAP@0.5:0.95} \\
        \midrule
        \textbf{Overall} & \textbf{0.736} & \textbf{0.740} & \textbf{0.745} & \textbf{0.448} \\
        Heavy-Vehicle & 0.913 & 0.978 & 0.981 & 0.763 \\
        Light-Vehicle & 0.892 & 0.951 & 0.961 & 0.649 \\
        Pedestrian & 0.822 & 0.915 & 0.918 & 0.522 \\
        Crack & 0.576 & 0.484 & 0.505 & 0.240 \\
        Crack-Severe & 0.548 & 0.503 & 0.493 & 0.273 \\
        Pothole & 0.597 & 0.440 & 0.468 & 0.198 \\
        Speed-Bump & 0.804 & 0.908 & 0.885 & 0.487 \\
        \bottomrule
    \end{tabular}
\end{table}

\subsubsection*{Video Processing Validation}

Sample processed video demonstrating real-world performance:

\begin{itemize}
    \item \textbf{Test Video}: 3-minute dashcam footage (1080p @ 30 FPS)
    \item \textbf{Processing Time}: 7.2 minutes (0.42x real-time on CPU-only mode)
    \item \textbf{Detections}: 47 potholes, 23 cracks, 8 speed bumps
    \item \textbf{False Positives}: 3 instances (shadows misclassified as cracks)
    \item \textbf{Video Reference}: \url{https://www.youtube.com/watch?v=W_SFcuZuRBE}
\end{itemize}

\subsection{System Integration Testing}

\begin{table}[H]
    \centering
    \caption{Integration Test Scenarios}
    \label{tab:integration-tests}
    \begin{tabular}{p{0.25\textwidth}p{0.5\textwidth}p{0.15\textwidth}}
        \toprule
        \textbf{Test Scenario} & \textbf{Description} & \textbf{Status} \\
        \midrule
        Image Upload & Single image processed through Streamlit UI & Pass \\
        Video Upload & 5-minute video file processed and downloaded & Pass \\
        Live Camera & Real-time detection from USB webcam & Pass \\
        Model Switching & Toggle between M1/M2 during processing & Pass \\
        Confidence Tuning & Adjust threshold and verify detection count changes & Pass \\
        Session Persistence & Upload $\rightarrow$ reload page $\rightarrow$ download remains available & Pass \\
        Concurrent Users & Multiple Streamlit sessions (load test) & \textit{Not Tested} \\
        GPS Integration & Frame-to-GPS synchronization & \textit{Pending} \\
        \bottomrule
    \end{tabular}
\end{table}

\clearpage

\section{Performance Optimization}

\subsection{Inference Optimization Techniques}

\begin{itemize}
    \item \textbf{Batch Processing}: Process multiple frames simultaneously (GPU parallelism).
    \item \textbf{Frame Sampling}: Process every Nth frame for video (e.g., 1 FPS instead of 30 FPS).
    \item \textbf{Model Quantization}: Convert FP32 models to FP16 or INT8 for edge devices.
    \item \textbf{TensorRT}: Optimize YOLO models for NVIDIA GPUs (potential 3--5x speedup).
\end{itemize}

\subsection{Network Bandwidth Optimization}

For edge-to-cloud data transfer:

\begin{itemize}
    \item \textbf{Image Compression}: JPEG quality 85\% reduces file size by 60--70\%.
    \item \textbf{Selective Upload}: Only upload frames with detections above confidence threshold.
    \item \textbf{Metadata-Only Mode}: Upload bounding box coordinates + GPS, store images locally.
\end{itemize}

\section{Security Considerations}

Although the current repository implementation lacks production security features, the following measures should be implemented for deployment:

\begin{table}[H]
    \centering
    \caption{Security Recommendations}
    \label{tab:security}
    \begin{tabular}{p{0.25\textwidth}p{0.65\textwidth}}
        \toprule
        \textbf{Area} & \textbf{Recommendation} \\
        \midrule
        Authentication & Implement JWT-based API authentication for edge device uploads \\
        Data Encryption & Use HTTPS/TLS for all API communication; encrypt local SQLite database \\
        Model Security & Sign model files to prevent tampering; validate checksums before loading \\
        Input Validation & Sanitize file uploads; enforce max file size limits (100 MB) \\
        Access Control & Restrict dashboard access via OAuth2 or institutional SSO \\
        Audit Logging & Log all detection uploads, model changes, and user actions \\
        \bottomrule
    \end{tabular}
\end{table}

\section{Future Enhancements}

\subsection{Short-Term (3--6 Months)}

\begin{enumerate}
    \item Implement GPS hardware integration with timestamp synchronization.
    \item Add persistent PostgreSQL database with spatial indexing (PostGIS).
    \item Develop RESTful API using FastAPI or Flask for production deployment.
    \item Implement message queue (RabbitMQ) for asynchronous processing.
\end{enumerate}

\subsection{Long-Term (6--12 Months)}

\begin{enumerate}
    \item Train improved models on larger, more diverse datasets (weather conditions, night driving).
    \item Implement severity classification using depth estimation or multi-frame analysis.
    \item Develop mobile iOS/Android app for field data collection.
    \item Integrate with city GIS systems via standard APIs (WFS, WMS).
    \item Add predictive maintenance features (pothole growth tracking over time).
\end{enumerate}

\section{Conclusion}

The \texttt{road-anomaly-detection} repository provides a robust foundation for the ML Inference subsystem of the Pothole Detection System architecture. The YOLOv8-based detection pipeline achieves state-of-the-art performance (0.745 mAP@0.5) and supports flexible deployment scenarios ranging from desktop analysis to edge computing.

\medskip

The Streamlit web dashboard offers immediate usability for prototype validation, while the modular codebase facilitates integration with production backend infrastructure. With the addition of GPS hardware integration and a production-grade API layer, the system will fulfill all requirements defined in the architectural design document.

\medskip

Key strengths of the current implementation:

\begin{itemize}
    \item High detection accuracy on diverse road anomaly types.
    \item Real-time processing capability (30+ FPS with GPU acceleration).
    \item Flexible multi-model architecture for ensemble detection.
    \item User-friendly interface for rapid prototyping and validation.
\end{itemize}

\medskip

This technical document serves as a bridge between the high-level architectural design and the practical implementation, providing clear guidance for system deployment, testing, and future development.

\end{document}
